% RAFTcorr User Guide -- preamble (pdflatex, xelatex or lualatex).

\documentclass[11pt,letterpaper]{article}

% ---------------------------------------------------------------- fonts
\usepackage[T1]{fontenc}
\usepackage[utf8]{inputenc}
\usepackage{lmodern}
\usepackage{amssymb}
\usepackage{amsmath}
\usepackage{textcomp}
\usepackage{microtype}
\DisableLigatures[-]{encoding=T1,family=tt*}   % --roi stays two hyphens in typewriter text

% ---------------------------------------------------------------- layout
\usepackage[margin=1in,headheight=14pt]{geometry}
\usepackage{fancyhdr}
\pagestyle{fancy}
\fancyhf{}
\fancyhead[L]{\small RAFTcorr User Guide}
\fancyhead[R]{\small\thepage}
\renewcommand{\headrulewidth}{0.4pt}
\setlength{\parskip}{4pt}

% ---------------------------------------------------------------- graphics
\usepackage{graphicx}
\graphicspath{{figures/}}
\usepackage{float}
\usepackage[font=small,labelfont=bf,margin=0.5in]{caption}

% ---------------------------------------------------------------- color
\usepackage{xcolor}
\definecolor{raftaccent}{RGB}{79,70,229}     % the application's accent (light theme)
\definecolor{raftgreen}{RGB}{22,163,74}
\definecolor{raftred}{RGB}{220,38,38}
\definecolor{raftamber}{RGB}{180,83,9}
\definecolor{codebg}{RGB}{245,245,247}
\definecolor{keybg}{RGB}{243,244,246}

% ---------------------------------------------------------------- hyperlinks
\usepackage[colorlinks=true,linkcolor=raftaccent,urlcolor=raftaccent,
            citecolor=raftaccent]{hyperref}

% ---------------------------------------------------------------- callouts
\usepackage[most]{tcolorbox}
\newtcolorbox{tipbox}{
  colback=raftgreen!7!white, colframe=raftgreen, boxrule=0.5pt,
  arc=1pt, left=4pt, right=4pt, top=3pt, bottom=3pt,
  title={\textbf{Tip}}, fonttitle=\bfseries\small}
\newtcolorbox{warnbox}{
  colback=raftred!7!white, colframe=raftred, boxrule=0.5pt,
  arc=1pt, left=4pt, right=4pt, top=3pt, bottom=3pt,
  title={\textbf{Warning}}, fonttitle=\bfseries\small}
\newtcolorbox{notebox}{
  colback=raftaccent!7!white, colframe=raftaccent, boxrule=0.5pt,
  arc=1pt, left=4pt, right=4pt, top=3pt, bottom=3pt,
  title={\textbf{Note}}, fonttitle=\bfseries\small}

% ---------------------------------------------------------------- code
\usepackage{listings}
\lstset{
  basicstyle=\ttfamily\small,
  backgroundcolor=\color{codebg},
  frame=single, framerule=0pt,
  breaklines=true, breakindent=0pt,
  columns=fullflexible, keepspaces=true, showstringspaces=false,
  xleftmargin=4pt, xrightmargin=4pt,
  aboveskip=6pt, belowskip=6pt,
  literate={~}{{\textasciitilde}}1,
}

% ---------------------------------------------------------------- misc
\usepackage{multicol}
\usepackage{enumitem}
\setlist{topsep=3pt,itemsep=2pt,parsep=0pt}
\usepackage{booktabs}
\usepackage{tabularx}
\usepackage{array}

% ---------------------------------------------------------------- interface vocabulary
% \ui{Run DIC Analysis}     a label, button or field exactly as the application shows it
% \menu{File}{Save Project} a menu path (any number of levels: \menu{Help}{Models…} uses two)
% \kbd{Ctrl+S}              a key or a key combination
% \file{results.h5}         a file name, folder or command-line word
% \screenshot[width]{file}{caption}{label}   a figure from figures/
\newcommand{\ui}[1]{\textbf{#1}}
\newcommand{\menu}[2]{\textbf{#1}\,$\triangleright$\,\textbf{#2}}
\newcommand{\menuthree}[3]{\textbf{#1}\,$\triangleright$\,\textbf{#2}\,$\triangleright$\,\textbf{#3}}
\newcommand{\kbd}[1]{\tcbox[on line, boxsep=0pt, left=3pt, right=3pt, top=1.5pt, bottom=1.5pt,
  colback=keybg, colframe=black!30, boxrule=0.4pt, arc=2pt]{\footnotesize\ttfamily #1}}
\newcommand{\file}[1]{\texttt{#1}}
\newcommand{\screenshot}[4][\textwidth]{%
  \begin{figure}[htbp]
    \centering
    \includegraphics[width=#1]{#2}
    \caption{#3}
    \label{#4}
  \end{figure}}

% Unicode characters used in labels of the application
\DeclareUnicodeCharacter{2026}{\ldots}
\DeclareUnicodeCharacter{03B5}{\ensuremath{\varepsilon}}
\DeclareUnicodeCharacter{03C3}{\ensuremath{\sigma}}
\DeclareUnicodeCharacter{03C9}{\ensuremath{\omega}}
\DeclareUnicodeCharacter{0394}{\ensuremath{\Delta}}
\DeclareUnicodeCharacter{2081}{\ensuremath{_1}}
\DeclareUnicodeCharacter{2082}{\ensuremath{_2}}
\DeclareUnicodeCharacter{2080}{\ensuremath{_0}}
\DeclareUnicodeCharacter{2192}{\ensuremath{\rightarrow}}
\DeclareUnicodeCharacter{00D7}{\ensuremath{\times}}
\DeclareUnicodeCharacter{2264}{\ensuremath{\leq}}
\DeclareUnicodeCharacter{2265}{\ensuremath{\geq}}
